% !TeX program = xelatex
\documentclass[11pt]{report}

\RequirePackage[banglamainfont=Kalpurush, banglattfont=SiyamRupali, feature=2]{latexbangla_exp}

\usepackage{ragged2e}
\usepackage{blindtext}

\usepackage{epigraph}
\usepackage{amsmath,amsfonts,amsthm,amssymb,graphicx,multicol,wrapfig,relsize}
\usepackage{tikz}
\usepackage[most]{tcolorbox}
\usepackage{marginnote}
\usepackage[top=4cm, bottom=3.5cm, outer=2cm, inner=2cm, right=3.7cm, left=2.5cm, marginparwidth=2cm, marginparsep=2cm]{geometry} 
\setlength{\columnsep}{0.24in}
\usepackage{header}
\usepackage{hyperref}
\hypersetup{
    colorlinks=true,
    linkcolor=blue,
    filecolor=magenta,
    urlcolor=cyan}
\urlstyle{same}
\title{Diophantine Equations}
\author{মোঃ ফোয়াদ আল আলম}
\date{February 2023}
%--------------------------------------------------------------%
\begin{document}

\pdfpageheight=11.25in
\pdfpagewidth=8in

\maketitle

\tableofcontents

\chapter{সূচনা}

\section{ডায়োফ্যান্টাইন সমীকরণ কী জিনিস???}

আমাদের পাঠ্যবইয়ের অনেক অঙ্কে সমীকরণ সমাধান করতে বলা হয়। বীজগণিত শিক্ষার শুরু থেকেই তোমরা নানা ধরনের সমীকরণ সমাধান করে এসেছ। এদের বেশিরভাগ সমীকরণকেই আমরা \textcolor{blue}{বীজগাণিতিক} হিসেবে আখ্যায়িত করতে পারি। অর্থাৎ, এধরনের সমীকরণে আমরা সকল বাস্তব সংখ্যার জন্য সমাধান বের করি। কিন্তু যে সকল সমীকরণে শুধু পূর্ণসংখ্যা অথবা ধনাত্মক পূর্ণসংখ্যার জন্য সমাধান বের করা হয়, তাদেরকে ডায়োফ্যান্টাইন সমীকরণ বলা হয়। ডায়োফ্যান্টাইন শব্দটি এসেছে গ্রিক গণিতবিদ \href{https://en.wikipedia.org/wiki/Diophantus}{ডায়োফ্যান্টাস} এর নাম থেকে।\\

চল আমরা কিছু ডায়োফ্যান্টাইন সমীকরণের উদাহরণ এবং সেগুলো কীভাবে সমাধান করতে হয় সেটা দেখি।\\

\section{সহজ একটা সমীকরণ}

নিচের সমীকরণটি লক্ষ কর: $$x^2+y^2=2$$ এই সমীকরণের বাস্তব সংখ্যায় অসংখ্য সমাধান আছে। যদি তুমি $|x| \leq \sqrt{2}$ এর জন্য $x$ এর যেকোনো বাস্তব মান প্রদত্ত সমীকরণে বসাও তাহলে $y$ এর এমন বাস্তব মান পাবে যা সমীকরণটিকে সিদ্ধ করে।\\

এখন উপরের সমীকরণের পূর্ণসাংখ্যিক সমাধান নিয়ে চিন্তা কর। এই সমীকরণের কি অসংখ্য পূর্ণসাংখ্যিক সমাধান বিদ্যমান?\\

না! এই সমীকরণের ক্ষেত্রে $x, y$ এর মান শুধু $\pm 1$ হওয়া সম্ভব। লক্ষ কর, সকল পূর্ণসংখ্যার বর্গের সেট, $$\{0^2, (\pm 1)^2, (\pm 2)^2, (\pm 3)^2, \ldots\} = \{0, 1, 4, 9, \ldots\}$$ অর্থাৎ, $|x| \geq 2$ হলে $2=x^2+y^2 \geq x^2 \geq 4,$ একটি কন্ট্রাডিকশন! অর্থাৎ, $|x| \leq 1$ হতে হবে। অনুরূপভাবে, $|y| \leq 1.$ এখান থেকে পাওয়া যায়, $(x,y)=(\pm 1, \pm 1).$\\

এই আইডিয়া আরো জটিল বিভিন্ন সমস্যা সমাধানেও কাজে লাগে।\\

\section{পারলে ফ্যাক্টরাইজ কর!}

বিভিন্ন ধরনের ডায়োফ্যান্টাইন সমীকরণে চলক সম্বলিত একটি রাশিকে ফ্যাক্টরাইজ করার মাধ্যমে সমাধান বের করা যায়। নিচের উদাহরণটি লক্ষ কর।\\

\begin{tcolorbox}[colframe=black!80!green,title=উদাহরণ 1.3.1]
পূর্ণসংখ্যায় সমাধান বের কর $$(3x+2y)(3y+2x)=17.$$
\end{tcolorbox}

খেয়াল কর, $3x+2y$ ও $3y+2x$ উভয়ই পূর্ণসংখ্যা। এই দুইটি পূর্ণসংখ্যা গুণ করলে $17$ পাওয়া যায়। অর্থাৎ, এই পূর্ণসংখ্যা জোড়াকে অবশ্যই $(1, 17), (17, 1), (-1, -17), (-17, -1)$ এদের কোনো একটি হতে হবে।\\

\textbf{\textcolor{blue}{কেস ১:}} $$3x+2y=1$$ $$3y+2x=17$$

\textbf{\textcolor{blue}{কেস ২:}} $$3x+2y=17$$ $$3y+2x=1$$

\textbf{\textcolor{blue}{কেস ৩:}} $$3x+2y=-1$$ $$3y+2x=-17$$

\textbf{\textcolor{blue}{কেস ৪:}} $$3x+2y=-17$$ $$3y+2x=-1$$

প্রত্যেকটি কেস আলাদাভাবে বিবেচনা করে দেখা যায়, এই সমীকরণের কোনো পূর্ণসাংখ্যিক সমাধান নেই।\\

এই সমস্যায় আমরা বিভাজ্যতার আইডিয়া ব্যবহার করেছি। তাই এই চারটি কেসের বাইরে অন্য কোনো কেস থাকা সম্ভব নয়। এই সমস্যায় প্রদত্ত সমীকরণটি ফ্যাক্টরাইজড অবস্থায়ই ছিল। কিন্তু প্রদত্ত সমীকরণ ফ্যাক্টরাইজড আকারে না থাকলেও চলক সম্বলিত রাশিকে ফ্যাক্টরাইজ করার মাধ্যমে সমাধান নির্ণয় করা সম্ভব। এখন আরেকটি উদাহরণ লক্ষ কর।\\

\begin{tcolorbox}[colframe=black!80!green,title=উদাহরণ 1.3.2]
এমন সকল পূর্ণসংখ্যা $x, y$ নির্ণয় কর যেন $$x^2=12+y^2.$$
\end{tcolorbox}

প্রদত্ত সমীকরণ থেকে আমরা পাই, $$x^2-y^2=12$$ $$\implies (x+y)(x-y)=12$$
এখানে দুইটি পূর্ণসংখ্যার গুণফল $12.$ অতএব, এদেরকে অবশ্যই $\pm (1, 12), \pm (2, 6), \pm (3, 4)$ অথবা এদের Permutation হতে হবে (যদি $(3, 4)$ একটি সমাধান হয়, তবে এর Permutation $(4, 3)$-ও একটি সমাধান)। এখন লক্ষ কর, $x+y=a$ এবং $x-y=b$ হলে সমীকরণ দুটিকে যোগ করে পাওয়া যায়, $2x = a+b.$ অর্থাৎ, $a+b$ কে জোড় সংখ্যা হতে হবে। $(a,b)=(6,2),(2,6),(-6,-2),(-2,-6)$ এই চারটি কেসেই শুধু $a+b$ এর মান জোড় হয়।\\

উপরের চারটি কেস সমাধান করে পাই, $(x,y)=(4,2),(4,-2),(-4,2),(-4,-2).$\\

\section{Simon's Favourite Factorizing Trick}

বিভিন্ন সমীকরণ সমাধান করতে গেলে অনেক ক্ষেত্রে ফ্যাক্টরাইজ করার আইডিয়া তেমন সুস্পষ্ট না।\\

\begin{tcolorbox}[colframe=black!80!green,title=উদাহরণ 1.4.1]
এমন সকল পূর্ণসংখ্যা $x, y$ নির্ণয় কর যেন $$xy=x+y+3.$$
\end{tcolorbox}

প্রদত্ত সমীকরণকে এভাবে লেখা যায় $$xy-x-y=3.$$ এখন সমীকরণের বামপক্ষকে কোনোভাবে ফ্যাক্টরাইজ করতে পারলে ভাল হত না?\\

\begin{tcolorbox}[colframe=blue!80!blue]
যখন শুধুমাত্র $xy, x$ এবং $y$ পদযুক্ত রাশি দেখবে, তখন নিচের অভেদটি ব্যবহার করতে পারো: $$xy+ax+by+ab=(x+b)(y+a)$$ অর্থাৎ, যদি এমন একটি সমীকরণ পাও $$xy+ax+by=c$$ তাহলে সমীকরণের উভয়পক্ষের সাথে $ab$ যোগ করলে বামপক্ষকে ফ্যাক্টরাইজ করে ফেলা সম্ভব। $$xy+ax+by+ab=c+ab \implies (x+b)(y+a)=c+ab$$ এই ট্রিককে অনেক সময় Simon's Favourite Factorizing Trick বা সংক্ষেপে SFFT হিসেবে সম্বোধন করা হয়। 
\end{tcolorbox}

প্রদত্ত সমীকরণের উভয়পক্ষের সাথে $1$ যোগ করলে পাই, $$xy-x-y+1=4 \implies (x-1)(y-1) = 4$$ এখান থেকে তোমরা সমাধান বের করতে পারবে না?\\

\begin{tcolorbox}[colframe=black!80!green,title=উদাহরণ 1.4.2]
এমন সকল মৌলিক সংখ্যা $p, q, r$ নির্ণয় কর যেন $$ pqr = 41 (p+q+r). $$
\end{tcolorbox}

লক্ষ কর, $ 41 | pqr ,$ যেখানে $41$ নিজেই একটি মৌলিক সংখ্যা। যেহেতু $pqr$ তিনটি মৌলিক সংখ্যার গুণফল, সেহেতু $p, q$ অথবা $r$ এর যেকোনো একটিকে অবশ্যই $41$ হতে হবে। WLOG, $p$ এর মান $41$। প্রদত্ত সমীকরণে $p$ এর মান বসিয়ে পাই, $$qr = q+r+41 \implies qr-q-r=41$$
এখন SFFT ব্যবহার করে পাই, $$qr-q-r+1=42 \implies (q-1)(r-1)=42$$
এখান থেকে তোমরা সহজেই দেখাতে পারবে যে, $(q,r) \in \{(2, 43), (43, 2)\}.$\\

\section{বর্গ পূরণ কর!}

কোনো সমীকরণের চলক সম্বলিত রাশিকে কতিপয় বর্গের সমষ্টিরূপে লিখতে পারলে বর্গগুলোকে বাউন্ড করার মাধ্যমে সমাধান নির্ণয় করা সম্ভব। নিচের উদাহরণগুলো দেখলে বিষয়টা বুঝতে পারবে।\\

\begin{tcolorbox}[colframe=black!80!green,title=INMO]
এমন সকল অঋণাত্মক পূর্ণসংখ্যা $x, y$ নির্ণয় কর যেন $$(xy-7)^2=x^2+y^2.$$
\end{tcolorbox}

আমরা প্রদত্ত সমীকরণকে এভাবে লিখতে পারি $$(xy)^2-14xy+49=x^2+y^2$$ $$ \implies (xy)^2-12xy+49=x^2+2xy+y^2$$ $$\implies (xy)^2-12xy+36+13=(x+y)^2$$ $$\implies (xy-6)^2+13=(x+y)^2$$ $$\implies (x+y)^2-(xy-6)^2=13$$ $$\implies (x+y+xy-6)(x+y-xy+6)=13$$ অতএব, $(x+y+xy-6, x+y-xy+6) = (13,1), (1,13)$ হতে হবে। এখান থেকে তোমরা সহজেই বের করতে পারবে যে, $(x,y)=(0,7),(7,0),(3,4),(4,3).$\\

\begin{tcolorbox}[colframe=black!80!green,title=উদাহরণ 1.5.2]
এমন সকল পূর্ণসংখ্যা $x, y$ নির্ণয় কর যেন $$x^2+y^2=x+y+2.$$
\end{tcolorbox}

ধর আমাদের কাছে এই সমীকরণটি দেওয়া ছিল $$x^2+y^2=2x+2y+2.$$ তাহলে আমরা বর্গগুলো পূরণ করতে পারি এবং উপরের সমীকরণকে এভাবে লিখতে পারি $$(x-1)^2+(y-1)^2=4$$ এখান থেকে বলা যায় $(x-1, y-1)=(0,\pm 2), (\pm 2,0)$ হতে হবে। কিন্তু প্রদত্ত সমীকরণে বর্গ পূরণ করার আইডিয়া এতটা সুস্পষ্ট না।\\

আমাদের এখন একটু চালাকি করতে হবে। যখনই তোমার কাছে $x^2 \pm x$ থাকবে, সমীকরণের উভয়পক্ষকে $4$ দ্বারা গুণ করলে বর্গ পূরণ করা সম্ভব। অর্থাৎ, আমরা প্রদত্ত সমীকরণ থেকে পাই, $$4x^2 - 4x + 4y^2 - 4y = 8 \implies (2x-1)^2 + (2y-1)^2 = 10.$$ অতএব, $(2x-1,2y-1)= \pm (1,3), \pm (3,1).$ এখন সমাধানগুলো তোমরা বের করে ফেল।\\

\begin{tcolorbox}[colframe=black!80!green,title=উদাহরণ 1.5.3]
এমন সকল পূর্ণসংখ্যা $x, y$ নির্ণয় কর যেন $$x^3+y^3=(x+y)^2.$$
\end{tcolorbox}

প্রদত্ত সমীকরণের বামপক্ষকে আমরা সহজেই সূত্র প্রয়োগ করে ফ্যাক্টরাইজ করতে পারি। $$(x+y)(x^2-xy+y^2)=(x+y)^2$$ এখন দুইটি কেস হওয়া সম্ভব।\\

\textbf{\textcolor{blue}{কেস ১:}} $x+y=0.$\\

এক্ষেত্রে আমরা সমাধান হিসেবে পাই $(x,y) = (a, -a)$ যেখানে $a \in \mathbb{Z}.$\\

\textbf{\textcolor{blue}{কেস ২:}} $x+y \not = 0.$\\

উপরোক্ত সমীকরণের উভয়পক্ষকে $(x+y)$ দ্বারা ভাগ করে পাই, $$x^2-xy+y^2=x+y$$ $$\implies x^2-xy+y^2-x-y=0$$ এখানে কি কোনোভাবে বর্গ পূরণ করা যাবে?\\

হ্যাঁ, যাবে! সমীকরণের উভয়পক্ষকে $2$ দ্বারা গুণ করে পাই $$2x^2-2xy+2y^2-2x-2y=0$$ $$\implies (x-y)^2+(x-1)^2+(y-1)^2=2$$ এখান থেকে এই কেসের সমাধানগুলো বের করা তোমাদের কাজ।\\

\section{The Legendary চিপা ট্রিক?!}

সকল পূর্ণসংখ্যার বর্গের সেট নিয়ে আবার চিন্তা কর। $$\{0^2, (\pm 1)^2, (\pm 2)^2, (\pm 3)^2, \ldots\} = \{0, 1, 4, 9, \ldots\}$$ ধর, এদের মাঝে দুইটি ক্রমিক পূর্ণবর্গসংখ্যা $m^2$ এবং $(m+1)^2$। তোমাকে যদি এমন কোনো পূর্ণসংখ্যা $n$ দেওয়া হয় যেন $m^2 < n < (m+1)^2$, সেক্ষেত্রে $n$ কি কখনো পূর্ণবর্গসংখ্যা হওয়া সম্ভব?\\

না! দুইটি ক্রমিক পূর্ণবর্গসংখ্যার মাঝে অন্য কোনো পূর্ণবর্গসংখ্যা থাকা সম্ভব নয়। আবার যদি কোনো পূর্ণবর্গসংখ্যা $n$ এর জন্য $(m-1)^2 < n < (m+1)^2$ হয়, তবে $n$ কে অবশ্যই $m^2$ এর সমান হতে হবে। এই ফ্যাক্টগুলো খুব অবভিয়াস হলেও বিভিন্ন সমস্যা সমাধানের ক্ষেত্রে এটি কার্যকরী ভূমিকা রাখতে পারে।\\

এই ট্রিকটি বাংলাদেশ জাতীয় গণিত ক্যাম্পে \textcolor{blue}{চিপা}(?!) ট্রিক হিসেবে পরিচিত। কারণ এই ট্রিক ব্যবহার করে আমরা $X < Y < Z$ এরকম কিছু দেখানোর চেষ্টা করি, যা দেখে মনে হয় $Y$ সংখ্যাটি $X$ এবং $Z$ এর চিপার মধ্যে রয়েছে। নিচের উদাহরণগুলো দেখলে এই ট্রিকের ব্যবহার ভালো বুঝতে পারবে।\\

\begin{tcolorbox}[colframe=black!80!green,title=APMO 1999]
এমন সকল ধনাত্মক পূর্ণসংখ্যা $a, b$ নির্ণয় কর যেন $a^2+4b$ এবং $b^2+4a$ উভয়ই পূর্ণবর্গসংখ্যা।
\end{tcolorbox}

$a^2+4b$ এবং $b^2+4a$ রাশি দুইটি পরস্পর প্রতিসম, মানে রাশি দুইটিতে $a$ এর স্থানে $b$ এবং $b$ এর স্থানে $a$ বসালে রাশি দুইটি একটি অপরটিতে পরিণত হয়। তাই WLOG, $a \leq b$ ধরে নেওয়া যায়।\\

আমরা পাই, $$b^2 < b^2+4a \leq b^2+4b < b^2+4b+4 = (b+2)^2.$$ অর্থাৎ, $b^2+4a$ রাশিটি $b^2$ এবং $(b+2)^2$ রাশি দুইটির চিপায় আছে। অতএব, $b^2+4a$ রাশিটি পূর্ণবর্গ হতে হলে $b^2+4a = (b+1)^2$ হতে হবে! কিন্তু, $$b^2+4a = (b+1)^2 \implies 4a = 2b + 1$$ একটি কন্ট্রাডিকশন! (কেন?)\\

আইডিয়াটা অনেক চমৎকার, তাই না? এই ট্রিকটি নিচের মত সমস্যা সমাধানে খুবই কার্যকরী।\\

\begin{tcolorbox}[colframe=blue!80!blue]
এমন সকল পূর্ণসংখ্যা $x$ নির্ণয় কর যেন $f(x)$ কোনো পূর্ণসংখ্যার পারফেক্ট $n$-তম ঘাত।
\end{tcolorbox}

এরকম সমস্যার ক্ষেত্রে যদি কোনোভাবে দেখানো যায় যে, কোনো পূর্ণসংখ্যা $y$-এর জন্য $y^n<x<(y+1)^n$, তবে এমন কোনো $x$ থাকা সম্ভব নয়।\\

চলো আরেকটি উদাহরণ দেখি।\\

\begin{tcolorbox}[colframe=black!80!green,title=উদাহরণ 1.6.2]
এমন সকল ধনাত্মক পূর্ণসংখ্যা $n$ নির্ণয় কর যেন $n^2-19n+89$ একটি পূর্ণবর্গসংখ্যা।
\end{tcolorbox}

যেহেতু $n$ এর সহগ $19$ একটি বিজোড় সংখ্যা, সুতরাং একে দুইটি জোড় সংখ্যার চিপায় ফেলা সম্ভব। অর্থাৎ, ধ্রুবপদ বাদ দিয়ে চিন্তা করলে, $$n^2-20n<n^2-19n<n^2-18n$$ ধ্রুবপদগুলো আমাদের কাঙ্খিত অসমতার ক্ষেত্রে সামঞ্জস্যতা বজায় রাখে না। তবে মূল রাশিটিতে $n^2$ সম্বলিত পদ বিদ্যমান থাকায় $n$ এর মান একটি নির্দিষ্ট সংখ্যা থেকে বড় হলে আমাদের কাঙ্খিত অসমতা সবসময় সত্য হওয়া উচিত।\\

ঠিক তাই হবে! একটু হিসাব কষলে দেখা যায়, $n > 11$ এর জন্য $n^2-20n+100 < n^2-19n+89$। আবার $n>8$ এর জন্য $n^2-19n+89 < n^2-18n+81$। অসমতা দুটিকে একত্রে লেখা যায়, $n>11$ এর জন্য, $$n^2-20n+100 < n^2-19n+89 < n^2-18n+81$$ $$\implies (n-10)^2 < n^2-19n+89 < (n-9)^2$$ অর্থাৎ, $n>11$ হলে $n^2-19n+89$ রাশিটি পূর্ণবর্গ হওয়া সম্ভব নয়। তার মানে $n \leq 11$ হতে হবে।\\

এখন $n$ এর মান প্রদত্ত রাশিতে বসিয়ে চেক করলে দেখা যায়, শুধু $n=8,11$ এর জন্য প্রদত্ত রাশিটি পূর্ণবর্গ হওয়া সম্ভব।\\

এই সমস্যায় আমরা চিপা ট্রিক ব্যবহার করে চলকের জন্য বাউন্ড বের করেছি। এরপর সীমিতসংখ্যক কেসের জন্য মান বসিয়ে চেক করে সমাধানগুলো বের করা সম্ভব।\\

\newpage

\section{নিশ্চায়কের ব্যবহার}

আমাদের স্কুলের পাঠ্যবইয়ের বীজগণিত অংশে দ্বিঘাত সমীকরণের সমাধান নির্ণয় করা শেখানো হয়। চলক $x$ এর জন্য একটি আদর্শ দ্বিঘাত সমীকরণ $ax^2+bx+c=0$ এর সমাধানদ্বয়, $$x= \frac{-b \pm \sqrt{b^2-4ac}}{2a}$$ এখানে, রুটের ভিতরের $b^2-4ac$ রাশিটির উপর সমাধানের প্রকৃতি নির্ভর করে। $b^2-4ac$ রাশিকে সমীকরণের নিশ্চায়ক বা ইংরেজিতে \href{https://en.wikipedia.org/wiki/Discriminant}{Discriminant} হিসেবে সম্বোধন করা হয়। কোনো সমীকরণের নিশ্চায়ককে $D$ বা $\Delta$ দ্বারা চিহ্নিত করা হয়। বিভিন্ন ডায়োফ্যান্টাইন সমীকরণের সমাধানে নিশ্চায়ক সম্পর্কিত এই ফ্যাক্টটি ব্যবহার করা যায়।\\

\begin{tcolorbox}[colframe=blue!80!blue]
$ax^2+bx+c=0$ সমীকরণের সমাধানদ্বয় মূলদ (আমাদের ব্যবহারের ক্ষেত্রে পূর্ণসংখ্যা) হতে হলে নিশ্চায়ক $\Delta = b^2-4ac$ রাশিকে অবশ্যই পূর্ণবর্গ হতে হবে। অর্থাৎ, $$\Delta = b^2-4ac=k^2$$ যেখানে $k \in \mathbb{Z}$.
\end{tcolorbox}

আমরা আবার উদাহরণ 1.6.2 এ ফিরে যাই।\\

\begin{tcolorbox}[colframe=black!80!green,title=উদাহরণ 1.6.2]
এমন সকল ধনাত্মক পূর্ণসংখ্যা $n$ নির্ণয় কর যেন $n^2-19n+89$ একটি পূর্ণবর্গসংখ্যা।
\end{tcolorbox}

আমরা এই সমস্যায় এর আগে চিপা ট্রিক ব্যবহার করে সমাধান বের করেছি। এখন আমরা নিশ্চায়কের বৈশিষ্ট্য ব্যবহার করে সমস্যা সমাধান করার চেষ্টা করব।\\

সমীকরণ বানানোর জন্য প্রদত্ত তথ্য অনুযায়ী $n^2-19n+89=z^2$ লেখা যায়, যেখানে $z \in \mathbb{Z}$. এখন কোনোভাবে $z$ এর মান বের করতে পারলেই দ্বিঘাত সমীকরণ সমাধান করে সমাধানগুলো বের করা সম্ভব।\\

সমীকরণটিকে নিচের মত করে সাজিয়ে লেখা যায়। $$n^2-19n+(89-z^2)=0$$ যা চলক $n$ এর জন্য একটি দ্বিঘাত সমীকরণ নির্দেশ করে। সমীকরণের নিশ্চায়ক $\Delta = (\pm 19)^2 - 4 \cdot 1 \cdot (89-z^2)$ নিশ্চয়ই একটি পূর্ণবর্গসংখ্যা! (কেন?!) অর্থাৎ, $$\Delta = (\pm 19)^2 - 4 \cdot 1 \cdot (89-z^2) = a^2$$ $$\implies a^2 - 4z^2=5$$ $$\implies (a+2z)(a-2z) = 5$$ অতএব, $(a+2z,a-2z) = (\pm 1,\pm 5),(\pm 5,\pm 1)$ হতে হবে। একটু হিসাব করে দেখ, $z = \pm 1$ হতে হবে। এখন, $$n^2-19n+89=z^2=1$$ $$\implies n^2-19n+88=0$$ $$\implies (n-8)(n-11)=0$$ অর্থাৎ, $n=8, 11$-ই নির্ণেয় সমাধান, যা আমরা এর আগেই একবার বের করেছিলাম!\\

\section{অনুশীলনী}

\begin{flushleft}

\textbf{\textcolor{blue}{সমস্যা ১.}} এমন সকল ধনাত্মক পূর্ণসংখ্যা $x, y, z$ নির্ণয় কর যেন $$xyz+x+y+z=2+xy+yz+zx.$$\\

\textbf{\textcolor{blue}{সমস্যা ২.}} এমন সকল ধনাত্মক পূর্ণসংখ্যা $x, y, z$ নির্ণয় কর যেন $$x^3+y^3+z^3-3xyz=p$$ যেখানে $p$ একটি মৌলিক সংখ্যা।\\

\textbf{\textcolor{blue}{সমস্যা ৩.}} ধনাত্মক পূর্ণসংখ্যা $n$ এর সর্বনিম্ন কোন মানের জন্য নিচের সমীকরণে $(x,y)$ এর সর্বমোট $69$টি সমাধান বিদ্যমান? $$\frac{1}{x}+\frac{1}{y}=\frac{1}{n}$$\\

\textbf{\textcolor{blue}{সমস্যা ৪.}} এমন সকল ধনাত্মক পূর্ণসংখ্যা $x, y, z$ নির্ণয় কর যেন $$x^2y+y^2z+z^2x=3xyz.$$\\

\textbf{\textcolor{blue}{সমস্যা ৫.}} এমন সকল ধনাত্মক পূর্ণসংখ্যা $x, y, z, w$ নির্ণয় কর যেন $$x^2+y^2+z^2+2xy+2y(z-1)+2x(z+1)=w^2.$$\\

\textbf{\textcolor{blue}{সমস্যা ৬.}} এমন সকল ধনাত্মক পূর্ণসংখ্যা $x, y$ নির্ণয় কর যেন $$x^3+(x+1)^3+(x+2)^3+\ldots+(x+7)^3=y^3.$$\\

\textbf{\textcolor{blue}{সমস্যা ৭.}} এমন সকল ধনাত্মক পূর্ণসংখ্যা $x, y$ নির্ণয় কর যেন $$(x^2+y)(x+y^2)=(xy)^3.$$\\

\end{flushleft}

\chapter{মডুলার অ্যারিথমেটিকের ব্যবহার}

\section{সমতাকে কনগ্রুয়েন্সে পরিণত কর!}

বিভিন্ন ধরনের অলিম্পিয়াড প্রবলেমে সাধারণ কোনো পদ্ধতি বা ট্রিক ব্যবহার করে সমাধান নির্ণয় করা সম্ভব নয়। এসব ক্ষেত্রে মডুলার অ্যারিথমেটিক ব্যবহার করার মাধ্যমে সমাধান বের করতে হয়। প্রদত্ত সমীকরণকে কোনো একটি নির্দিষ্ট সংখ্যা $n$ এর মডুলোতে নিলে সমতাকে কনগ্রুয়েন্সে পরিণত করা যায়।\\

\begin{tcolorbox}[colframe=blue!80!blue]
\begin{center}
    বামপক্ষ = ডানপক্ষ $\implies$ বামপক্ষ $\equiv$ ডানপক্ষ $\pmod{n}$
\end{center} 
\end{tcolorbox}

এভাবে বিভিন্ন সমস্যায় সমীকরণের দুইপাশে mod $n$ নিলে কোনো কন্ট্রাডিকশনে উপনীত হওয়া সম্ভব। আবার অনেক সময় ঐ সমীকরণের বিভিন্ন চলকের ব্যাপারে নতুন তথ্য বের করা যায়, যা সমাধানের পরবর্তী ধাপে কাজে লাগানো যেতে পারে। এ ধরনের সমীকরণ সমাধানের ক্ষেত্রে $n$ এর কার্যকরী মান বের করতে পারাটাই সবচেয়ে কঠিন ধাপ। নিচের উদাহরণগুলো দেখলে বিষয়টি আরো ভালোভাবে বুঝতে পারবে।\\

\begin{tcolorbox}[colframe=black!80!green,title=উদাহরণ 2.1.1]
এমন সকল পূর্ণসংখ্যা $x, y$ নির্ণয় কর যেন $$y^2 = x^5 + 7.$$
\end{tcolorbox}

এখানে, সমীকরণের একটি চলকের ঘাত $5$ বিদ্যমান। এমন সমীকরণ কোনো সাধারণ পদ্ধতি দিয়ে সমাধান করা সাধারণত সম্ভব নয়।\\

সমীকরণটিকে mod $11$ এ নিলে পাই, $$y^2 \equiv x^5 + 7 \pmod{11}$$.
লক্ষ কর, mod $11$ এ $x^5 \in \{-1, 0, 1\}.$ অর্থাৎ, $y^2 \equiv x^5 +7 \equiv 6, 7, 8 \pmod{11}$ হতে হবে। কিন্তু একটু হিসাব করে দেখ, $y^2 \in \{0, 1, 3, 4, 5, 9\} \pmod{11},$ একটি কন্ট্রাডিকশন!\\

এই আপাত জটিল সমস্যার সমাধানটি খুবই ছোট এবং সুন্দর, এই ব্যাপারে তোমাদের মনে কোনো সন্দেহ থাকার কথা নয়। কিন্তু অনেকের মনেই প্রশ্ন জাগতে পারে, কারো মাথায় mod $11$ নেওয়ার চিন্তা কীভাবে আসবে?\\

এজন্য তোমাকে প্রচুর অনুশীলন করে ইনটুইশন বিল্ডআপ করতে হবে। এধরনের সমস্যা সমাধানে মডুলার অ্যারিথমেটিকের কিছু বৈশিষ্ট্য নিয়ে চিন্তা করলে কোন সংখ্যার mod নেওয়া উচিত তার জন্য ভাল ক্যান্ডিডেট বাছাই করা যায়।\\

মডুলার অ্যারিথমেটিক ব্যবহার করে কন্ট্রাডিকশন দেখানোর ক্ষেত্রে যেকোনো চলক সম্বলিত পদের সম্ভাব্য যত কমসংখ্যক রেসিডিউ পাওয়া সম্ভব, বামপক্ষ এবং ডানপক্ষের সম্ভাব্য সকল রেসিডিউয়ের সেট নিশ্ছেদ হিসেবে পাওয়ার সম্ভাবনা তত বৃদ্ধি পায়।\\

আগের সমস্যাটি নিয়ে আবার চিন্তা কর। আমাদের লক্ষ এমন কোনো $n$ এর mod নেওয়া যেন $x^5$ সম্বলিত পদটির খুব কম সংখ্যক সম্ভাব্য রেসিডিউ থাকে।\\

মডুলার অ্যারিথমেটিক শিখতে যেয়ে তোমরা নিশ্চয়ই \href{https://en.wikipedia.org/wiki/Fermat%27s_little_theorem}{Fermat's Little Theorem} এর ব্যাপারে জেনেছ? Fermat's Little Theorem
অনুযায়ী, যেকোনো মৌলিক সংখ্যা $p$ এর জন্য

\begin{center}
$a^{p-1} \equiv 1 \pmod{p};$ যখন $p \nmid a$\\
$a^{p-1} \equiv 0 \pmod{p};$ যখন $p \mid a$
\end{center}

এখান থেকে নিচের সিদ্ধান্তে উপনীত হওয়া সম্ভব।\\

\begin{tcolorbox}[colframe=blue!80!blue]

যেকোনো মৌলিক সংখ্যা $p$ এর জন্য $a^{\frac{p-1}{2}} \in \{0, \pm 1\} \pmod{p}.$

\end{tcolorbox}

$p=11$ হলে $x^5 \equiv x^{\frac{11-1}{2}} \in \{-1, 0, 1\} \pmod{11}$. এটাই আমাদের mod $11$ নিয়ে চেক করার অন্যতম প্রধান কারণ।\\

$k=\frac{p-1}{2}$ ধরলে $p=2k+1$, অর্থাৎ কোনো সমীকরণে যদি $x^k$ সম্বলিত কোনো পদ থাকে, তবে $2k+1$ মৌলিক সংখ্যা হয় কিনা চেক করে দেখতে পারো। যদি $2k+1$ মৌলিক সংখ্যা হয়ে থাকে, তবে পুরো সমীকরণকে $2k+1$ এর মডুলোতে নিলে কাজে লাগতে পারে।\\

এখন আমরা আরেকটি উদাহরণ দেখব।\\

\begin{tcolorbox}[colframe=black!80!green,title=উদাহরণ 2.1.2]
এমন সকল পূর্ণসংখ্যা $x, y$ নির্ণয় কর যেন $$x^2 - y! = 2001.$$
\end{tcolorbox}

লক্ষ কর, সমীকরণের বামপক্ষে ফ্যাক্টরিয়ালযুক্ত পদ বিদ্যমান। পুরো সমীকরণকে কোনো সংখ্যার মডুলোতে নিলে ফ্যাক্টরিয়ালযুক্ত পদ অপসারণের জন্য নিচের ফ্যাক্টটির সুবিধা নেওয়া যায়।\\

\begin{tcolorbox}[colframe=blue!80!blue]

কোনো স্বাভাবিক সংখ্যা $n$ এর জন্য, $y \geq n$ হলে, $n \mid y!.$

\end{tcolorbox}

চলো সমীকরণটিকে $7$ এর মডুলোতে নিই। সমীকরণটিকে এভাবে লেখা যায়, $$x^2 - y! \equiv -1 \pmod{7}$$ এখন আমরা দুইটি কেস নিয়ে চিন্তা করতে পারি।\\

\textbf{\textcolor{blue}{কেস ১:}} $y \geq 7$\\

এক্ষেত্রে $7 \mid y! \implies y! \equiv 0 \pmod{7}.$ তাহলে নিচের কনগ্রুয়েন্সকে অবশ্যই সত্য হতে হবে। $$x^2 \equiv -1 \pmod{7}$$ কিন্তু একটু হিসাব করে দেখ, $x^2 \in \{0,1,2,4\} \pmod{7}$, একটি কন্ট্রাডিকশন!\\

\textbf{\textcolor{blue}{কেস ২:}} $0 \leq y \leq 6$\\

এক্ষেত্রে $y$ এর মান $x^2=2001+y!$ সমীকরণে বসিয়ে চেক করে পাই, শুধু $y=4$ এর জন্য $x^2 = 2001 + 4! = 2025 = (\pm 45)^2.$ অর্থাৎ, $(x,y) = (45, 4), (-45, 4)$ সমীকরণকে সিদ্ধ করে।\\

এই সমস্যায় আমরা মডুলার অ্যারিথমেটিক ব্যবহার করে কীভাবে একটি চলককে (এখানে $y$) বাউন্ড করা যায় তা দেখলাম। এরপর সেই চলকের সীমিত সংখ্যক কেসের জন্য মান বসিয়ে চেক করে সমাধান বের করা সম্ভব।\\

\newpage

\section{কিসের মডুলোতে নিব?}

সমীকরণে কোনো চলকের ঘাত বড় হলে সমাধানের ক্ষেত্রে মডুলার অ্যারিথমেটিকের ব্যবহার অনেক কার্যকরী ভূমিকা রাখতে পারে। শুধু কন্ট্রাডিকশন দেখানোর মাধ্যমে নয়, একটি নির্দিষ্ট চলকের ব্যাপারে গুরুত্বপূর্ণ তথ্য বের করার ক্ষেত্রেও মডুলার অ্যারিথমেটিক ব্যবহার করা যায়। তবে কোন ঘাতের জন্য কিসের mod নেওয়া বেশিরভাগ ক্ষেত্রে কার্যকরী, সে ব্যাপারে ধারণা থাকলে কোনো সমস্যার সমাধান বের করার কাজটা সহজ হয়ে যায়। নিচে এমনই কিছু কাজের সম্পর্ক দেখানো হলো।\\

\begin{tcolorbox}[colframe=blue!80!blue]

প্রদত্ত সমীকরণের কোনো চলকের ঘাত $2$ হলে পুরো সমীকরণকে $3, 4, 5, 8$ ইত্যাদি সংখ্যার মডুলোতে নিলে গুরুত্বপূর্ণ তথ্য পাওয়া যেতে পারে।\\
$$a^2 \equiv \{0,1\} \pmod{3}$$
$$a^2 \equiv \{0,1\} \pmod{4}$$
$$a^2 \equiv \{0, \pm 1\} \pmod{5}$$
$$a^2 \equiv \{0,1,4\} \pmod{8}$$
বিশেষ করে $2$ এর ঘাতের (যেমন- $4, 8, 16$ ইত্যাদি) সাপেক্ষে mod নেওয়া অনেক কার্যকরী।\\

প্রদত্ত সমীকরণের কোনো চলকের ঘাত $3$ হলে পুরো সমীকরণকে $7$ অথবা $9$ এর মডুলোতে নিলে গুরুত্বপূর্ণ তথ্য পাওয়া যেতে পারে।\\
$$a^3 \equiv \{0, \pm 1\} \pmod{7}$$
$$a^3 \equiv \{0, \pm 1\} \pmod{9}$$

\end{tcolorbox}

\newpage

\section{সূচকীয় সমীকরণের পূর্ণসাংখ্যিক সমাধান}

এতক্ষণ পর্যন্ত আমরা যত সমীকরণ দেখেছি তাদের সবগুলোই বীজগাণিতিক রাশি সম্বলিত সমীকরণ ছিল। সূচকীয় সমীকরণ সমাধানের ক্ষেত্রে মডুলার অ্যারিথমেটিক ব্যবহার করা ছাড়াও অনেক সময় কিছু অ্যাডভান্সড টেকনিক ব্যবহার করার প্রয়োজন হয়। তবে এখন আমরা তুলনামূলক সহজ একটি সূচকীয় সমীকরণের উদাহরণ দেখব।\\

\begin{tcolorbox}[colframe=black!80!green,title=উদাহরণ 2.3.1]
এমন সকল ধনাত্মক পূর্ণসংখ্যা $x, y$ নির্ণয় কর যেন $$3^x-2^y=7.$$
\end{tcolorbox}

খেয়াল কর, $y \geq 3$ হলে, $8 \mid 2^y$. এখন পুরো সমীকরণকে mod $8$ এ নিলে পাওয়া যায়, $3^x \equiv 7 \pmod{8}$. কিন্তু $3^x \in \{1,3\} \pmod{8}$, একটি কন্ট্রাডিকশন!\\

$y \leq 2$ এর জন্য কোনো সমাধান পাওয়া যায় কিনা সেটা তোমরা চেক কর।\\

\section{ফ্যাক্টর-মডুলার ডুয়ো!}

অনেক সময় ফ্যাক্টরাইজ করার পাশাপাশি মডুলার অ্যারিথমেটিক ব্যবহার করার মাধ্যমে বিভিন্ন সমস্যা সমাধান করা যায়। নিচের উদাহরণটি দেখলেই বিষয়টি বুঝতে পারবে।\\

\begin{tcolorbox}[colframe=black!80!green,title=উদাহরণ 2.4.1]
এমন সকল পূর্ণসংখ্যা $x, y$ নির্ণয় কর যেন $$y^2=x^3+23.$$
\end{tcolorbox}

এ সমস্যা সমাধানে \href{https://web.evanchen.cc/handouts/ORPR/ORPR.pdf}{অর্ডার} দ্বারা প্রমাণ করা যায় এমন একটি বিখ্যাত ফ্যাক্ট ব্যবহার করতে হবে।\\

\begin{tcolorbox}[colframe=blue!80!blue]

যদি কোনো মৌলিক সংখ্যা $p \equiv 3 \pmod{4}$ হয়, তবে $-1$ কখনোই $p$ এর \href{https://en.wikipedia.org/wiki/Quadratic_residue}{কুয়াড্রাটিক রেসিডিউ} নয়; বিপরীতভাবেও এটি সত্য। $$p \equiv 3 \pmod{4} \iff n^2+1 \not \equiv 0 \pmod{p}$$

\end{tcolorbox}

যদি $y$ বিজোড় সংখ্যা হয়, তবে $y^2 \equiv 1 \pmod{4}.$ এখন প্রদত্ত সমীকরণকে mod $4$ এ নিয়ে পাই, $$x^3 \equiv 1-23 \equiv 2 \pmod{4}$$ 
যা অসম্ভব। তার মানে $y$ কে জোড় সংখ্যা হতে হবে। $y$ জোড় সংখ্যা হলে, $$x^3 \equiv 0-23 \equiv 1 \pmod{4}$$ $$\implies x \equiv 1 \pmod{4}$$

এখন, প্রদত্ত সমীকরণকে এভাবে লেখা যায়, $$y^2+4 = x^3+27$$ $$\implies 4(\frac{y^2}{4}+1)  = x^3+3^3$$ $$\implies 4((\frac{y}{2})^2+1) = (x+3)(x^2-3x+9)$$ খেয়াল কর, $x^2-3x+9 \equiv 3 \pmod{4}$. তার মানে এমন একটি মৌলিক সংখ্যা $p \equiv 3 \pmod{4}$ আছে যেন $p \mid x^2-3x+9$. (কেন?) তার মানে, $$p \mid (\frac{y}{2})^2+1$$ অর্থাৎ, $-1$, $p$ এর একটি কুয়াড্রাটিক রেসিডিউ, একটি কন্ট্রাডিকশন!

\newpage

\section{পীথাগোরিয়ান ট্রিপল}

তিনটি ধনাত্মক পূর্ণসংখ্যার ট্রিপল $(a, b, c)$ কে \href{https://en.wikipedia.org/wiki/Pythagorean_triple}{পীথাগোরিয়ান ট্রিপল} বলা হবে যদি এবং কেবল যদি $$a^2+b^2=c^2.$$ পীথাগোরিয়ান ট্রিপল নামটি অবশ্যই \href{https://en.wikipedia.org/wiki/Pythagorean_theorem}{পীথাগোরাসের উপপাদ্য} থেকে অনুপ্রাণিত।\\

দুনিয়ার সকল পীথাগোরিয়ান ট্রিপলকে তালিকাভুক্ত করা সম্ভব! খেয়াল কর, gcd$(a,b,c)=k$ হলে, আমরা সমীকরণের উভয়পক্ষকে $k^2$ দিয়ে ভাগ করতে পারি। তাই ধরে নিই, gcd$(a,b,c)=1$. এমন সকল পীথাগোরিয়ান ট্রিপলকে আমরা \textcolor{blue}{প্রিমিটিভ} বলে সম্বোধন করব।\\

\begin{tcolorbox}[colframe=blue!80!blue]

$(a,b,c)$ একটি প্রিমিটিভ পীথাগোরিয়ান ট্রিপল হলে, এমন দুটি পূর্ণসংখ্যা $m$ ও $n$ বিদ্যমান যেন $$(a,b,c) = (2mn, m^2-n^2, m^2+n^2).$$

\end{tcolorbox}

$a,b,c$ যদি এই আকারের হয়, তবে সহজেই চেক করা যায় যে $a^2+b^2=c^2$ হবে। কিন্তু সবচেয়ে মজার বিষয় হলো এটি বিপরীতভাবেও সত্য। তোমরা এটি প্রমাণ করার চেষ্টা করতে পারো।\\

এখন আমরা এই তথ্যটি ব্যবহার করে কীভাবে একটি সমস্যা সমাধান করা যায় সেটা দেখব।\\

\begin{tcolorbox}[colframe=black!80!green,title= INMO 2018]
$\triangle ABC$ এমন একটি ত্রিভুজ যেন এর তিনটি বাহুর দৈর্ঘ্যই পূর্ণসংখ্যা এবং কোনো দুইটি বাহুর দৈর্ঘ্য পরস্পর সমান নয়। $D$ ও $E$ যথাক্রমে $BC$ ও $CA$ বাহুর মধ্যবিন্দু এবং $\triangle ABC$ এর ভরকেন্দ্র $G$। $D, C, G$ ও $E$ একই বৃত্তের উপর অবস্থিত হলে $\triangle ABC$ এর সম্ভাব্য সর্বনিম্ন পরিসীমা নির্ণয় কর।
\end{tcolorbox}

কিছু জ্যামিতিক দক্ষতা প্রয়োগ করলে পাওয়া যায় যে, প্রদত্ত সমস্যায় $(a,b,c) \in \mathbb{N}^3$ বের করতে বলা হয়েছে যেন $a^2+b^2=2c^2$ এবং $a+b+c$ এর মান সর্বনিম্ন হয়। (এটুকু নিজেরা দেখানোর চেষ্টা কর)\\

উপরের সমীকরণকে এভাবে লেখা যায় $$(2c)^2 = 2 \cdot 2c^2 = 2(a^2+b^2) = (a+b)^2 + (a-b)^2$$ অর্থাৎ, $(a-b,a+b,2c)$ একটি পীথাগোরিয়ান ট্রিপল গঠন করে।\\

সবচেয়ে ছোট পীথাগোরিয়ান ট্রিপল হলো $(3,4,5)$, কিন্তু এটি কাজ করবে না। (কেন?) এর পরবর্তী পীথাগোরিয়ান ট্রিপল হলো $(6,8,10)$। তবে সর্বনিম্ন কার্যকরী পীথাগোরিয়ান ট্রিপল পেতে হলে প্রিমিটিভ পীথাগোরিয়ান ট্রিপল চেক করতে হবে। (কেন?) তাই আমরা $(5,12,13)$ নিয়ে চেষ্টা করি। পূর্ণসংখ্যায় পরিণত করতে স্কেল করার পর পাওয়া যায় যে, $(a,b,c) = (17,7,13)$ অর্থাৎ, $17+7+13=37$-ই আমাদের কাঙ্খিত পরিসীমা।\\

\section{Infinite Descent}

বিভিন্ন ডায়োফ্যান্টাইন সমীকরণ সমাধানের ক্ষেত্রে \href{https://en.wikipedia.org/wiki/Proof_by_infinite_descent}{Infinite Descent} একটি অতি গুরুত্বপূর্ণ ট্রিক। এই ট্রিক ব্যবহারের ক্ষেত্রে আমরা ধরে নিই যে, সমীকরণের কিছু সমাধানের অস্তিত্ব বিদ্যমান এবং আমরা সবচেয়ে \textcolor{blue}{ছোট} সমাধান নিয়ে কাজ করি। এরপর আমরা দেখাই যে, এর চেয়েও ছোট সমাধান বিদ্যমান, যা একটি কন্ট্রাডিকশন প্রদান করে। একটি উদাহরণ দেখলে বিষয়টি বুঝতে পারবে।\\

\begin{tcolorbox}[colframe=black!80!green,title=উদাহরণ 2.6.1]
এমন সকল ধনাত্মক পূর্ণসংখ্যা $x,y,z$ নির্ণয় কর যেন $$x^3+2y^3=4z^3.$$
\end{tcolorbox}

খেয়াল কর, $2 \mid x^3 \implies 2 \mid x$. এখানে, $x=2x'$ হলে প্রদত্ত সমীকরণকে এভাবে লেখা যায়, $$8x'^3+2y^3=4z^3$$ $$\implies 4x'^3+y^3=2z^3$$ অনুরূপভাবে, $2 \mid y^3 \implies 2 \mid y$. $y=2y'$ হলে প্রদত্ত সমীকরণকে এভাবে লেখা যায়, $$4x'^3+8y'^3=2z^3$$ $$\implies 2x'^3+4y'^3=z^3$$ অনুরূপভাবে, $2 \mid z^3 \implies 2 \mid z$. $z=2z'$ হলে প্রদত্ত সমীকরণকে এভাবে লেখা যায়, $$2x'^3+4y'^3=8z'^3$$ $$\implies x'^3+2y'^3=4z'^3$$

তার মানে $(x,y,z)$ একটি সমাধান হলে $(x',y',z')$-ও একটি সমাধান হবে। \href{https://en.wikipedia.org/wiki/Well-ordering_principle}{Well-ordering Principle} অনুসারে এমন একটি সমাধান $(a,b,c)$ বিদ্যমান যেন $a+b+c$ এর মান সর্বনিম্ন হয়। আবার, $(a,b,c)$ একটি সমাধান হলে $(\frac{a}{2},\frac{b}{2},\frac{c}{2})$-ও সমীকরণের একটি সমাধান। কিন্তু $$a+b+c > \frac{a}{2}+\frac{b}{2}+\frac{c}{2}$$ একটি কন্ট্রাডিকশন! অতএব, প্রদত্ত সমীকরণের পূর্ণসংখ্যায় কোনো সমাধান নেই।\\

\section{অনুশীলনী}

\begin{flushleft} 

\textbf{\textcolor{blue}{সমস্যা ১.}} পীথাগোরিয়ান ট্রিপলের সাধারণ সমাধান সম্পর্কিত উপপাদ্যটি প্রতিপাদন কর (সেকশন 2.5 দ্রষ্টব্য)।\\

\textbf{\textcolor{blue}{সমস্যা ২.}} এমন সকল পূর্ণসংখ্যা $x, y$ নির্ণয় কর যেন $$x^3+y^4=7.$$\\

\textbf{\textcolor{blue}{সমস্যা ৩.}} প্রমাণ কর, $m,n$ স্বাভাবিক সংখ্যা হলে $$3^m+3^n+1$$ কখনোই পূর্ণবর্গসংখ্যা নয়।\\

\textbf{\textcolor{blue}{সমস্যা ৪.}} এমন সকল ধনাত্মক পূর্ণসংখ্যা $m,n$ নির্ণয় কর যেন $$1!+2!+3!+\ldots+n!=m^2.$$\\

\textbf{\textcolor{blue}{সমস্যা ৫.}} দেখাও যে, নিচের সমীকরণের পূর্ণসংখ্যায় কোনো সমাধান নেই। $$a^3+(a+1)^3+(a+2)^3+\ldots+(a+6)^3 = b^4+(b+1)^4$$\\

\textbf{\textcolor{blue}{সমস্যা ৬.}} এমন সকল ধনাত্মক পূর্ণসংখ্যা $m,n$ নির্ণয় কর যেন $$3^m-7^n=2.$$\\

\textbf{\textcolor{blue}{সমস্যা ৭.}} দেখাও যে, $a^2+b^2=3c^2$ সমীকরণের ধনাত্মক পূর্ণসংখ্যায় কোনো সমাধান নেই।\\

\textbf{\textcolor{blue}{সমস্যা ৮.}} এমন সকল পূর্ণসংখ্যা $a,b,c,n$ নির্ণয় কর যেন $$6(6a^2+3b^2+c^2)=5n^2.$$\\

\textbf{\textcolor{blue}{সমস্যা ৯.}} এমন সকল পূর্ণসংখ্যা $x,y$ নির্ণয় কর যেন $$y^3=3x^2+3x+7.$$\\

\textbf{\textcolor{blue}{সমস্যা ১০.}} এমন সকল ট্রিপল $(m,p,q)$ নির্ণয় কর যেন $$2^mp^2+1=q^7$$ যেখানে $m \in \mathbb{N}$ এবং $p,q$ মৌলিক সংখ্যা।\\

\textbf{\textcolor{blue}{সমস্যা ১১.}} এমন সকল ট্রিপল $(x,y,p)$ নির্ণয় কর যেন $$x^5+x^4+1=p^y$$ যেখানে, $x,y \in \mathbb{N}$ এবং $p$ মৌলিক সংখ্যা।\\

\textbf{\textcolor{blue}{সমস্যা ১২.}} এমন সকল ধনাত্মক পূর্ণসংখ্যা $a,b,c$ নির্ণয় কর যেন $$a!b!=a!+b!+c!.$$\\

\textbf{\textcolor{blue}{সমস্যা ১৩.}} এমন সকল অঋণাত্মক পূর্ণসংখ্যা $x, y, z$ নির্ণয় কর যেন $$2^x + 3^y = z^2.$$\\

\textbf{\textcolor{blue}{সমস্যা ১৪.}} এমন সকল অঋণাত্মক পূর্ণসংখ্যা $x,y,z,w$ নির্ণয় কর যেন $$2^x3^y-5^z7^w=1.$$\\

\end{flushleft}


\chapter{অলিম্পিয়াডের কিছু সমস্যা}

এ অধ্যায়ে আমরা দেশ-বিদেশের বিভিন্ন অলিম্পিয়াডের কিছু ডায়োফ্যান্টাইন সমীকরণ এবং সেগুলো সমাধানে এখন পর্যন্ত শেখা বিভিন্ন টেকনিক ও ট্রিকস এর প্রয়োগ দেখব।\\

\section{$\Delta$ কোথায়?}

এখন আমরা একটি বিখ্যাত গণিত অলিম্পিয়াডের একটি সমস্যা দেখব। সমস্যাটি ১৯৮৭ সালে ইউএসএ-এর জাতীয় গণিত অলিম্পিয়াডের ১ নম্বর সমস্যা ছিল।\\

\begin{tcolorbox}[colframe=black!80!green,title=USAMO 1987]
এমন সকল পূর্ণসংখ্যা $x,y$ নির্ণয় কর যেন $$(x^2+y)(x+y^2)=(x-y)^3.$$
\end{tcolorbox}

প্রদত্ত সমীকরণকে নিচের মত করে সাজিয়ে পাই, $$x^2y^2+3x^2y-3xy^2+xy+2y^3=0$$ $$\implies y(x^2y+3x^2-3xy+x+2y^2)=0$$

$y=0$ হলে, $(x,y) = (a,0)$ প্রদত্ত সমীকরণের অসংখ্য সমাধান, যেখানে $a \in \mathbb{Z}$।\\

$y \not = 0$ হলে, উভয়পক্ষকে $y$ দ্বারা ভাগ করে পাই, $$x^2y+3x^2-3xy+x+2y^2=0$$ $$\implies 2y^2+(x^2-3x)y+(3x^2+x)=0$$

চলক $y$ এর সাপেক্ষে দ্বিঘাত সমীকরণটির নিশ্চায়ক, $$\Delta  = (x^2-3x)^2 - 4 \cdot 2 \cdot (3x^2+x)=x(x-8)(x+1)^2$$ এই রাশিটি অবশ্যই পূর্ণবর্গ হতে হবে। তার মানে $x(x-8)$ রাশিটিও পূর্ণবর্গ হতে হবে।\\

কোনো ধনাত্মক পূর্ণসংখ্যা $a,b$ এর জন্য $a^2-b^2=16$ হলে, $(a+b)(a-b)=16$. এখান থেকে পাওয়া যায়, $(a,b)=(4,0),(5,3)$.\\

যেহেতু $(x-4)^2-x(x-8)=16$, উপরের সমীকরণের সাথে তুলনা করে পাই, $x(x-8) \in \{0,9\}$. এখন, চলক $x$ এর জন্য সমাধান করে পাই, $x \in \{-1,0,8,9\}$.\\

চলক $y$ এর দ্বিঘাত সমীকরণটিতে $x$ এর মান বসিয়ে পাই, $(x,y)=(-1,-1),(8,-10),(9,-6),(9,-21).$\\

\section{পূর্ণবর্গ?}

এখন আমরা আরেকটি জাতীয় গণিত অলিম্পিয়াডের একটি সমস্যা দেখব। সমস্যাটি ১৯৯৭ সালে ইরানের জাতীয় গণিত অলিম্পিয়াডের ১ নম্বর সমস্যা ছিল।\\

\begin{tcolorbox}[colframe=black!80!green,title=Iran 2nd Round 1997]
ধনাত্মক পূর্ণসংখ্যা $x,y$ এর জন্য $3x^2+x=4y^2+y$. প্রমাণ কর, $x-y$ একটি পূর্ণবর্গসংখ্যা।
\end{tcolorbox}

লক্ষ কর, সমীকরণের বামপক্ষে $3$টি $x^3$ পদ থাকলেও ডানপক্ষে $4$টি $y^3$ পদ বিদ্যমান। যেহেতু $x-y$ রাশিটির বৈশিষ্ট্য বের করতে হবে, তাই প্রদত্ত সমীকরণে কোনোভাবে এই রাশিটিকে আনতে হবে। সমীকরণটিকে সাজিয়ে লিখলে পাই, $$3x^2+x=4y^2+y \iff y^2 = (x-y)(3x+3y+1) \iff x^2 = (x-y)(4x+4y+1).$$ দ্বিতীয় এবং তৃতীয় সমীকরণদ্বয় গুণ করে পাই, $$(xy)^2 = (x-y)^2(3x+3y+1)(4x+4y+1)$$ অর্থাৎ, $(3x+3y+1)(4x+4y+1)$ রাশিটি পূর্ণবর্গ। কিন্তু $$\text{gcd}(3x+3y+1,4x+4y+1) = \text{gcd}(3x+3y+1,x+y) = \text{gcd}(1,x+y)=1.$$ সুতরাং, $3x+3y+1$ এবং $4x+4y+1$ রাশিদ্বয় উভয়ই পূর্ণবর্গ হতে হবে। (কেন?) এখন, $y^2 = (x-y)(3x+3y+1)$ সমীকরণ থেকে বলা যায় $x-y$ রাশিটি পূর্ণবর্গ।\\

\section{ঘন ঘন ঘন!}

এখন আমরা কিছুটা পুরোনো একটি আইএমও শর্টলিস্টেড প্রবলেম দেখব। এই সমস্যাটি ২০০২ সালের আইএমও শর্টলিস্টের সংখ্যাতত্ত্ব অংশের ১ নম্বর সমস্যা ছিল।\\

\begin{tcolorbox}[colframe=black!80!green,title=IMO SL 2002]
ধনাত্মক পূর্ণসংখ্যা $t$ এর মান সর্বনিম্ন কত হলে এমন কিছু পূর্ণসংখ্যা $x_1, x_2, \ldots, x_t$ পাওয়া যাবে যেন $$x_1^3+x_2^3+\ldots+x_t^3=2002^{2002} ?$$
\end{tcolorbox}

প্রদত্ত সমীকরণে $x^3$ জাতীয় পদ দেখে কী করা উচিত বলে মনে হয়? হ্যাঁ, $7$ অথবা $9$ এর মডুলোতে নেওয়া উচিত! যেহেতু, $2002^{2002} \equiv 0 \pmod{7}$, তাই mod $7$ নিলে খুব একটা লাভ হবে বলে মনে হয় না। কিন্তু, $$2002^{2002} \equiv 4^{2002} \equiv (4^{333})^6 \cdot 4^4 \equiv 4 \pmod{9}$$

যেহেতু, $x^3 \in \{-1,0,1\} \pmod{9}$, সুতরাং সমীকরণের উভয়পক্ষ mod $9$ এ কনগ্রুয়েন্ট হতে হলে বামপক্ষে কমপক্ষে $4$টি পদ লাগবে ($1+1+1+1 \equiv 4 \pmod{9}$)। অর্থাৎ, $t \geq 4$ হতে হবে।\\

এখানেই কি আমাদের সমাধান শেষ? $t$ এর সর্বনিম্ন মান $4$ হবে এটা কি বলে দেওয়া যায়?\\

না! এধরনের সমস্যার সমাধানে দুইটি অংশ থাকে। প্রথমে চলকের উপর একটি বাউন্ড সৃষ্টি করে দেখাতে হয়, এরপর ঐ বাউন্ডই যে যথেষ্ট সেটা কন্সট্রাকশনের মাধ্যমে দেখাতে হয়। এই সমস্যার ক্ষেত্রে আমাদের $t=4$ এর জন্য নির্দিষ্ট কিছু পূর্ণসংখ্যা $x_1,x_2,x_3,x_4$ বের করে দেখাতে হবে যেন $x_1^3+x_2^3+x_3^3+x_4^3 = 2002^{2002}$ হয়।\\

তবে কন্সট্রাকশনের অংশটা দেখানো এতটা কঠিন নয়। $$2002^{2002} = 2002 \cdot (2002^{667})^3$$ $$= (10^3+10^3+1^3+1^3) \cdot (2002^{667})^3$$ $$= (10 \cdot 2002^{667})^3 + (10 \cdot 2002^{667})^3 + (2002^{667})^3 + (2002^{667})^3.$$

\newpage

\section{অনুশীলনী}

\begin{flushleft}

\textbf{\textcolor{blue}{সমস্যা ১.}} \textcolor{magenta}{(IMO 2006)} এমন সকল পূর্ণসংখ্যা $x,y$ নির্ণয় কর যেন $$1+2^x+2^{2x+1}=y^2.$$\\

\textbf{\textcolor{blue}{সমস্যা ২.}} \textcolor{magenta}{(USAJMO 2013)} এমন সকল পূর্ণসংখ্যা $a,b$ নির্ণয় কর যেন $a^5b+3$ এবং $ab^5+3$ উভয়ই পূর্ণঘনসংখ্যা।\\

\textbf{\textcolor{blue}{সমস্যা ৩.}} \textcolor{magenta}{(USAMO 2005)} এমন সকল পূর্ণসংখ্যা $x,y,z$ নির্ণয় কর যেন $$x^6+x^3+x^3y+y=147^{157}$$ $$x^3+x^3y+y^2+y+z^9=157^{147}.$$\\

\textbf{\textcolor{blue}{সমস্যা ৪.}} \textcolor{magenta}{(IMO SL 2010)} ধনাত্মক পূর্ণসংখ্যা $n$ এর সর্বনিম্ন কোন মানের জন্য পরস্পর ভিন্ন $n$টি ধনাত্মক পূর্ণসংখ্যার একটি সেট $\{s_1,s_2,\ldots,s_n\}$ বিদ্যমান যেন $$(1-\frac{1}{s_1})(1-\frac{1}{s_2}) \cdots (1-\frac{1}{s_n})= \frac{51}{2010} ?$$\\

\textbf{\textcolor{blue}{সমস্যা ৫.}} \textcolor{magenta}{(USAMO 2015)} এমন সকল ধনাত্মক পূর্ণসংখ্যা $x,y$ নির্ণয় কর যেন $$x^2+xy+y^2=(\frac{x+y}{3}+1)^2.$$\\

\textbf{\textcolor{blue}{সমস্যা ৬.}} \textcolor{magenta}{(IMO SL 2012)} এমন সকল পূর্ণসংখ্যার ট্রিপল $(x,y,z)$ নির্ণয় কর যেন $x \leq y \leq z$ এবং $$x^3(y^3+z^3)=2012(xyz+2).$$

\end{flushleft}

\begin{thebibliography}{9}

\bibitem{book} Aditya Khurmi, \emph{\href{https://drive.google.com/file/d/1BcJTLjQaelZ4w_70oHKyImC2I8zLfyrt/view}{Modern Olympiad Number Theory}}

\bibitem{handout} Atonu Roy Chowdhury, \emph{\href{https://atonurc.github.io/assets/dioph.pdf}{Diophantine Equations}}, BdMO National Camp 2021

\bibitem{handout} Navneel Singhal, \emph{\href{https://ghoshadi.files.wordpress.com/2020/05/diophantine-equations-note-1.pdf}{Diophantine Equations}}

\bibitem{handout} David Arthur, \emph{\href{https://drive.google.com/file/d/18l41QOLZ3M93z3lD7FKyE_CsXGsNbIvP/view}{Number Theory Tips and Tricks}}, Winter Camp 2009

\bibitem{handout} James Rickards, \emph{\href{https://drive.google.com/file/d/1K2RVdWEfb_NJ5b7shZsNExH2rJLEx1RH/view}{Diophantine Equations}}, Winter Camp 2019

\bibitem{website} Art of Problem Solving, \emph{\href{https://artofproblemsolving.com/}{AoPS}}

\end{thebibliography}

\end{document}